% Ch9 WS1 -- hybridization and sigma/pi bonding. Blocks 1-2. CED content.
\documentclass{apchem-worksheet}

\apcunitword{Chapter}
\apcunit{9}
\apcunittitle{Covalent Bonding: Orbitals}
\apcdoctitle{Worksheet 1 \textbullet\ Hybridization and $\sigma$/$\pi$ Bonds}
\apcsubtitle{Zumdahl \S9.1 \textbullet\ count domains, not bonds}

\begin{document}
\wsheader[Follow the three-step order every time: Lewis structure $\to$
VSEPR domains $\to$ hybridization. Show the domain count with every
hybridization answer.]

\problem[]{Explain, in your own words, why hybridization is proposed at all.
Use methane as the example. \answerlines[3]{Carbon's ground-state
configuration ($2s^2 2p^2$) predicts two bonds at 90$^\circ$, but methane
has four identical bonds at 109.5$^\circ$. Combining the $2s$ and three
$2p$ orbitals into four equivalent sp$^3$ orbitals produces exactly the
observed geometry --- the model is built to fit the molecule, not the
isolated atom.}}

\problem[]{Give the domain count and hybridization of the underlined atom:}
\begin{parts}
  \item C in \ce{CCl4}: \answerblank[4.5cm]{4 domains, sp$^3$}
  \item N in \ce{NH3}: \answerblank[4.5cm]{4 domains, sp$^3$}
  \item C in \ce{CO2}: \answerblank[4.5cm]{2 domains, sp}
  \item B in \ce{BCl3}: \answerblank[4.5cm]{3 domains, sp$^2$}
  \item O in \ce{H2O}: \answerblank[4.5cm]{4 domains, sp$^3$}
  \item S in \ce{SO2}: \answerblank[4.5cm]{3 domains, sp$^2$}
\end{parts}

\problem[]{A student assigns sp hybridization to the oxygen in water,
reasoning that oxygen forms two bonds. Identify the error and give the
correct answer. \answerlines[3]{Hybridization is set by electron
\emph{domains}, not bonds. Oxygen has two bonding pairs plus two lone pairs
--- four domains --- so it is sp$^3$, with a bent molecular shape and an
angle near 104.5$^\circ$.}}

\problem[]{Distinguish $\sigma$ and $\pi$ bonds:}
\begin{parts}
  \item How do the orbitals overlap in each?
        \answerlines[2]{$\sigma$: head-on along the internuclear axis.
        $\pi$: side-by-side, giving lobes above and below the axis.}
  \item Which type is formed from unhybridized $p$ orbitals?
        \answerblank[2.4cm]{$\pi$}
  \item Which type permits free rotation? \answerblank[2.4cm]{$\sigma$}
  \item Why is the first bond between two atoms always $\sigma$?
        \answerlines[2]{Head-on overlap is more effective than side-by-side,
        so it is energetically favoured and forms first.}
\end{parts}

\problem[]{Count the $\sigma$ and $\pi$ bonds in each molecule:}
\begin{parts}
  \item \ce{CH4}: \answerblank[4cm]{4 $\sigma$, 0 $\pi$}
  \item \ce{H2CO}: \answerblank[4cm]{3 $\sigma$, 1 $\pi$}
  \item \ce{C2H2}: \answerblank[4cm]{3 $\sigma$, 2 $\pi$}
  \item \ce{CO2}: \answerblank[4cm]{2 $\sigma$, 2 $\pi$}
  \item \ce{N2}: \answerblank[4cm]{1 $\sigma$, 2 $\pi$}
\end{parts}

\problem[]{Restricted rotation.}
\begin{parts}
  \item Explain why rotation about a C=C double bond is blocked but
        rotation about a C--C single bond is free.
        \answerlines[3]{A $\sigma$ bond is cylindrically symmetric about the
        axis, so rotation changes nothing. A $\pi$ bond requires its two $p$
        orbitals to stay parallel; rotating 90$^\circ$ would set them
        perpendicular and destroy the overlap, which means breaking the
        bond.}
  \item \ce{C2H2Cl2} with one Cl on each carbon exists as two separable
        compounds when the carbons are joined by a double bond, but not
        when joined by a single bond. Explain.
        \answerlines[3]{With a double bond the $\pi$ system locks the
        geometry, so \emph{cis} (both Cl on one side) and \emph{trans}
        (opposite sides) are distinct, isolable substances. With a single
        bond the two arrangements interconvert freely and are the same
        compound.}
\end{parts}

\problem[]{For acrylonitrile, \ce{H2C=CH-C#N}:}
\begin{parts}
  \item Hybridization of each carbon, left to right:
        \answerblank[5cm]{sp$^2$, sp$^2$, sp}
  \item Total $\sigma$ bonds: \answerblank[2cm]{6}
  \item Total $\pi$ bonds: \answerblank[2cm]{3}
  \item Approximate H--C--H angle:
        \answerblank[2.6cm]{$\approx120^\circ$}
\end{parts}

\begin{spiralstrip}{Chapters 7--8 \textbullet\ spiral}
  \item Configuration of \ce{Fe^2+}: \answerblank[3.2cm]{$[\ce{Ar}]3d^6$}
  \item $\Delta H$ from bond energies:
        \answerblank[4cm]{broken $-$ formed}
  \item Higher lattice energy: \ce{NaF} or \ce{MgO}?
        \answerblank[2.4cm]{\ce{MgO}}
  \item PES peak position represents: \answerblank[3.4cm]{binding energy}
  \item Is \ce{CO2} polar? \answerblank[2cm]{no}
\end{spiralstrip}

\end{document}
